\section{司法与仲裁}
\begin{enumerate}
  \item 仲裁裁决的执行。
        \begin{enumerate}
          \item 由被执行人住所地和被执行财产所在地的中院。\zd{看见仲裁找中院，除了国内仲裁的保全和国内仲裁的证据保全找基层法院}。
          \item 一个申请执行，一个申请撤销仲裁裁决 \yjt 中止执行。
        \end{enumerate}
  \item 司法对仲裁的监督。
        \begin{enumerate}
          \item 撤销仲裁裁决。
                \begin{enumerate}
                  \item 收到裁决之日起 6 个月内。
                  \item \zd{仲裁协议无效法定情形} 6 个 3 大类，非常重要。
                        \begin{enumerate}
                          \item 没有仲裁协议。
                          \item 仲裁裁决的事项不属于仲裁协议范围 \hz 仲裁委员会无权仲裁。
                          \item 仲裁庭的组成 \hz 仲裁程序违反法定程序的 \yjt \zd{程序错}。
                          \item 仲裁裁决所依据的证据伪造。
                          \item 对方隐瞒了足以影响公正裁决的证据的 \yjt \zd{伪造隐瞒证据}。
                          \item \zd{仲裁员徇私舞弊}。
                        \end{enumerate}
                  \item 要撤销仲裁裁决太难了，难到什么程度呢你当一辈子律师。
                        \begin{enumerate}
                          \item \xjt 两个说的是无效，但无效要在庭前提出，当时没提异议，后面都不准提了。
                          \item 仲裁基本上就没有非常严格的程序。
                          \item 法院对撤销仲裁裁决非常慎重。
                          \item 就是为了保证仲裁的效率，一裁终局；会影响仲裁公信力；仲裁制度的发展；增加法院负担。
                        \end{enumerate}
                  \item 2 个月之内组成合议庭审查。
                        \begin{enumerate}
                          \item 没有法定情形 \yjt 裁定驳回撤销申请。
                          \item 发现法定理由 \yjt 裁定撤销仲裁裁决。
                          \item 伪造隐瞒证据 \yjt 通知仲裁庭\zd{重新仲裁}，仲裁庭自行决定，如果重新仲裁不需要重组仲裁庭。
                        \end{enumerate}
                \end{enumerate}
          \item 仲裁裁决的不予执行同跟撤销仲裁裁决。
          \item 仲裁司法审查案件的报核程序。有些事儿是需要由法院来审查的。
                \begin{enumerate}
                  \item 确认仲裁协议效力、撤销仲裁裁决、不予执行仲裁裁决。
                  \item 否认仲裁（纠纷怎么解决，要么重新搞仲裁协议仲裁）必须走报核程序。
                  \item 由高院核准就行，涉外涉港澳台、公共利益最高院报核。
                \end{enumerate}
        \end{enumerate}
\end{enumerate}
